% Options for packages loaded elsewhere
\PassOptionsToPackage{unicode}{hyperref}
\PassOptionsToPackage{hyphens}{url}
\PassOptionsToPackage{dvipsnames,svgnames,x11names}{xcolor}
%
\documentclass[
  12pt,
  letterpaper,
]{book}

\usepackage{amsmath,amssymb}
\usepackage{iftex}
\ifPDFTeX
  \usepackage[T1]{fontenc}
  \usepackage[utf8]{inputenc}
  \usepackage{textcomp} % provide euro and other symbols
\else % if luatex or xetex
  \usepackage{unicode-math}
  \defaultfontfeatures{Scale=MatchLowercase}
  \defaultfontfeatures[\rmfamily]{Ligatures=TeX,Scale=1}
\fi
\usepackage{lmodern}
\ifPDFTeX\else  
    % xetex/luatex font selection
\fi
% Use upquote if available, for straight quotes in verbatim environments
\IfFileExists{upquote.sty}{\usepackage{upquote}}{}
\IfFileExists{microtype.sty}{% use microtype if available
  \usepackage[]{microtype}
  \UseMicrotypeSet[protrusion]{basicmath} % disable protrusion for tt fonts
}{}
\makeatletter
\@ifundefined{KOMAClassName}{% if non-KOMA class
  \IfFileExists{parskip.sty}{%
    \usepackage{parskip}
  }{% else
    \setlength{\parindent}{0pt}
    \setlength{\parskip}{6pt plus 2pt minus 1pt}}
}{% if KOMA class
  \KOMAoptions{parskip=half}}
\makeatother
\usepackage{xcolor}
\usepackage[margin=1in]{geometry}
\setlength{\emergencystretch}{3em} % prevent overfull lines
\setcounter{secnumdepth}{5}
% Make \paragraph and \subparagraph free-standing
\makeatletter
\ifx\paragraph\undefined\else
  \let\oldparagraph\paragraph
  \renewcommand{\paragraph}{
    \@ifstar
      \xxxParagraphStar
      \xxxParagraphNoStar
  }
  \newcommand{\xxxParagraphStar}[1]{\oldparagraph*{#1}\mbox{}}
  \newcommand{\xxxParagraphNoStar}[1]{\oldparagraph{#1}\mbox{}}
\fi
\ifx\subparagraph\undefined\else
  \let\oldsubparagraph\subparagraph
  \renewcommand{\subparagraph}{
    \@ifstar
      \xxxSubParagraphStar
      \xxxSubParagraphNoStar
  }
  \newcommand{\xxxSubParagraphStar}[1]{\oldsubparagraph*{#1}\mbox{}}
  \newcommand{\xxxSubParagraphNoStar}[1]{\oldsubparagraph{#1}\mbox{}}
\fi
\makeatother


\providecommand{\tightlist}{%
  \setlength{\itemsep}{0pt}\setlength{\parskip}{0pt}}\usepackage{longtable,booktabs,array}
\usepackage{calc} % for calculating minipage widths
% Correct order of tables after \paragraph or \subparagraph
\usepackage{etoolbox}
\makeatletter
\patchcmd\longtable{\par}{\if@noskipsec\mbox{}\fi\par}{}{}
\makeatother
% Allow footnotes in longtable head/foot
\IfFileExists{footnotehyper.sty}{\usepackage{footnotehyper}}{\usepackage{footnote}}
\makesavenoteenv{longtable}
\usepackage{graphicx}
\makeatletter
\def\maxwidth{\ifdim\Gin@nat@width>\linewidth\linewidth\else\Gin@nat@width\fi}
\def\maxheight{\ifdim\Gin@nat@height>\textheight\textheight\else\Gin@nat@height\fi}
\makeatother
% Scale images if necessary, so that they will not overflow the page
% margins by default, and it is still possible to overwrite the defaults
% using explicit options in \includegraphics[width, height, ...]{}
\setkeys{Gin}{width=\maxwidth,height=\maxheight,keepaspectratio}
% Set default figure placement to htbp
\makeatletter
\def\fps@figure{htbp}
\makeatother
% definitions for citeproc citations
\NewDocumentCommand\citeproctext{}{}
\NewDocumentCommand\citeproc{mm}{%
  \begingroup\def\citeproctext{#2}\cite{#1}\endgroup}
\makeatletter
 % allow citations to break across lines
 \let\@cite@ofmt\@firstofone
 % avoid brackets around text for \cite:
 \def\@biblabel#1{}
 \def\@cite#1#2{{#1\if@tempswa , #2\fi}}
\makeatother
\newlength{\cslhangindent}
\setlength{\cslhangindent}{1.5em}
\newlength{\csllabelwidth}
\setlength{\csllabelwidth}{3em}
\newenvironment{CSLReferences}[2] % #1 hanging-indent, #2 entry-spacing
 {\begin{list}{}{%
  \setlength{\itemindent}{0pt}
  \setlength{\leftmargin}{0pt}
  \setlength{\parsep}{0pt}
  % turn on hanging indent if param 1 is 1
  \ifodd #1
   \setlength{\leftmargin}{\cslhangindent}
   \setlength{\itemindent}{-1\cslhangindent}
  \fi
  % set entry spacing
  \setlength{\itemsep}{#2\baselineskip}}}
 {\end{list}}
\usepackage{calc}
\newcommand{\CSLBlock}[1]{\hfill\break\parbox[t]{\linewidth}{\strut\ignorespaces#1\strut}}
\newcommand{\CSLLeftMargin}[1]{\parbox[t]{\csllabelwidth}{\strut#1\strut}}
\newcommand{\CSLRightInline}[1]{\parbox[t]{\linewidth - \csllabelwidth}{\strut#1\strut}}
\newcommand{\CSLIndent}[1]{\hspace{\cslhangindent}#1}

% header.tex
\usepackage{xcolor}        % Para colores
\usepackage{polyglossia}   % Soporte de idiomas con XeLaTeX
\setdefaultlanguage{spanish}

% Definición de colores personalizados
\definecolor{fecundidad}{HTML}{FFC125}      % Amarillo
\definecolor{supervivencia}{HTML}{36648B}   % Azul
\definecolor{crecimiento}{HTML}{548B54}     % Verde
\definecolor{retroceso}{HTML}{7A378B}       % Violeta
\definecolor{permanencia}{HTML}{26466A}     % Azul oscuro
\makeatletter
\@ifpackageloaded{caption}{}{\usepackage{caption}}
\AtBeginDocument{%
\ifdefined\contentsname
  \renewcommand*\contentsname{Table of contents}
\else
  \newcommand\contentsname{Table of contents}
\fi
\ifdefined\listfigurename
  \renewcommand*\listfigurename{List of Figures}
\else
  \newcommand\listfigurename{List of Figures}
\fi
\ifdefined\listtablename
  \renewcommand*\listtablename{List of Tables}
\else
  \newcommand\listtablename{List of Tables}
\fi
\ifdefined\figurename
  \renewcommand*\figurename{Figure}
\else
  \newcommand\figurename{Figure}
\fi
\ifdefined\tablename
  \renewcommand*\tablename{Table}
\else
  \newcommand\tablename{Table}
\fi
}
\@ifpackageloaded{float}{}{\usepackage{float}}
\floatstyle{ruled}
\@ifundefined{c@chapter}{\newfloat{codelisting}{h}{lop}}{\newfloat{codelisting}{h}{lop}[chapter]}
\floatname{codelisting}{Listing}
\newcommand*\listoflistings{\listof{codelisting}{List of Listings}}
\makeatother
\makeatletter
\makeatother
\makeatletter
\@ifpackageloaded{caption}{}{\usepackage{caption}}
\@ifpackageloaded{subcaption}{}{\usepackage{subcaption}}
\makeatother

\ifLuaTeX
  \usepackage{selnolig}  % disable illegal ligatures
\fi
\usepackage{bookmark}

\IfFileExists{xurl.sty}{\usepackage{xurl}}{} % add URL line breaks if available
\urlstyle{same} % disable monospaced font for URLs
\hypersetup{
  pdftitle={Conclusión},
  colorlinks=true,
  linkcolor={Maroon},
  filecolor={Maroon},
  citecolor={Blue},
  urlcolor={Blue},
  pdfcreator={LaTeX via pandoc}}


\title{Conclusión}
\author{}
\date{}

\begin{document}
\frontmatter
\maketitle

\renewcommand*\contentsname{Table of contents}
{
\hypersetup{linkcolor=}
\setcounter{tocdepth}{2}
\tableofcontents
}

\mainmatter
\chapter{Por: Raymond L Tremblay, Demetria Mondragón y Aucencia
Emeterio-Lara}\label{por-raymond-l-tremblay-demetria-mondraguxf3n-y-aucencia-emeterio-lara}

La gama de acercamientos desarrollada en este libro permite ser
utilizado para evaluar una gran diversidad de hipótesis y temas. Pero,
sin duda lo más importante no son las técnicas matemáticas sino las
preguntas que se pueden responder con la información presentada. Las
hipótesis que se presentan desde el punto de vista de conservación,
ecología, historia de vida y evolución, son solo ejemplos de preguntas
que pueden ser respondidas con los modelos de dinámica poblacional. Gran
parte del libro está dedicado a resaltar los métodos cuantitativos y sus
limitaciones, reconociendo que para realizar análisis de dinámica
poblacional hay que entender todos los pasos y tener presente que, en
cualquiera de estos casos, existen errores cuantitativos y de
interpretación, que pueden invalidar el trabajo final. Los autores de
este libro reconocen que el estudio de la dinámica poblacional es
holístico y permite evaluar historias de vida complejas.

En esta sección presentamos algunas preguntas principales que pueden ser
abordadas en el futuro.

\chapter{Preguntas ecológica}\label{preguntas-ecoluxf3gica}

\begin{itemize}
\item
  Aunque la mayoría de las especies de orquídeas tienen solamente un
  polinizador, para aquella que tienen múltiples especies de
  polinizadores, no está claro si esto afecta la dinámica de sus
  poblaciones (Tremblay, 1992; Ackerman et al., 2025). En otras palabras
  ¿las especies con múltiples polinizadores tienen una mayor
  consistencia poblacional al compararlas con las especies con un solo
  polinizador? y de ser el caso ¿cuáles son los factores que lo están
  promoviendo? por ejemplo: ¿Cuando uno evalúa la diversidad de polínios
  sobre los polinizadores, esa diversidad pudiese tener repercusiones en
  la producción de frutos, calidad de las semillas y esto a su vez
  pudiera verse reflejado en la tasas de fecundidad y sobrevivencia de
  las plántulas, lo que sugiere un efecto sobre la dinámica de las
  poblaciones. Desafortunadamente se conoce muy poco sobre el efecto de
  la diversidad de padres sobre la calidad de las semillas (Tremblay,
  1994; Light \& MacConaill, 1998).
\item
  Típicamente los modelos están basados en las cifras de reclutamiento
  que se observa en el laboratorio (\emph{in vitro}), lo cual es una
  sub-estimación de la mortandad en el campo y una sobre-estimación del
  potencial de reclutamiento, ya que a dicho valor se le debiera
  incorporar eventos estocásticos como la proporción de semillas que
  llegan a un micronicho apropiado para la germinación, la proporción de
  semillas que germinan bajo las condiciones microambientales en las que
  se desarrollen, las perdidas por herbivoría o por desecación, etc. No
  está claro, si esta subestimación es relevante en la dinámica de las
  especies en ambientes naturales. La evaluación de la germinación de
  semillas en el campo: dado a su pequeño tamaño y la dificultad de
  detectarlas, así como los protocolos es un reto que necesita ser
  explorado con profundidad. Habrá necesidad de evaluar técnicas de
  germinación in vivo que permitan el monitoreo del reclutamiento a
  mediano plazo para tener cifras más exactas del número que integra la
  nueva generación de individuos (plántulas). En general se conoce muy
  poco sobre la proporción de las semillas producidas y los valores de
  germinación en el campo (Ackerman, Sabat \& Zimmerman, 1996; Ehrlén \&
  Eriksson, 2000; Batty et al., 2001; Wilcock \& Neiland, 2002).
\item
  El efecto de las variables abióticas sobre la historia de vida de las
  plantas comienzan a ser investigado (Morris et al., 2020), y su efecto
  sobre orquídeas es un tema que merece atención. Por ejemplo ¿Cuál es
  el efecto de los patrones de variación abiótica sobre la persistencia
  de las poblaciones y su dinámica a largo plazo?
\item
  Los efectos bióticos tales como la herbivoría, interacciones con
  polinizadores, las interacción con árboles forófitos (en el caso de
  orquídeas epífitas), la diversidad de micorrizas en la raíz de las
  orquídeas, etc pudieran influir sobre los parámetros de la modelos de
  proyección poblacional (PPM) , por lo que se necesita explorar a mayor
  profundidad dichos efectos.
\item
  Aunque los ecólogos que trabajan con orquídeas, reconocen que muchas
  especies se multiplican clonalmente, aparte de los trabajos de
  \emph{Cypripedium} (Garcı́a, Goni \& Guzmán, 2010) no es evidente su
  impacto sobre la persistencia y dinámica de estas poblaciones. Muchas
  otras especies de orquídeas tienen una historia de vida que incluye
  clonaje tal como \emph{Epidendrum}, \emph{Artorima} pero no se ha
  evaluado su impacto sobre la dinámica poblacional.
\item
  Los trabajos de orquídeas epífitas han determinado la dinámica
  poblacional con información de la orquídea, desligado del forófito,
  sin embargo considerando que no son plantas con crecimiento
  independiente del substrato, ¿de qué manera se puede incluir el efecto
  o importancia del forófito en los modelos demográficos. El único
  trabajo que evaluó el efecto del forófito sobre la dinámica
  poblacional de una orquídea epífita fue el de \emph{Oncidium
  brachyandrum} (Ramı́rez Martı́nez, 2021), en su tesis doctoral. El
  efecto de los diferentes hospederos de \emph{Quercus} es dramático,
  vea el capítulo de LTRE/ERTV.
\item
  ¿Cuál es la relación de la diversidad micorrízica, el ciclo de vida y
  la adecuación evolutiva? En otras palabras, ¿individuos con mayor
  diversidad micorrízica tienen una mayor adecuación evolutiva que
  individuos con menos diversidad hongos micorrízicos? O ¿los que tienen
  un simbionte específico tienen una mayor adecuación evolutiva,
  influenciando la edad a la primera reproducción o el esfuerzo
  reproductivo?
\end{itemize}

\chapter{Preguntas de historia de
vida}\label{preguntas-de-historia-de-vida}

\begin{itemize}
\item
  ¿Cuales son las diferencias y similitudes en la historia de vida de
  las orquídeas cuando se compara epífitas, terrestres, litófitas,
  templadas y tropicales?
\item
  Aunque hay poca diversidad de epífitas en áreas templadas, no son
  ausentes, por ejemplo \emph{Sarcochilus australis} es una epífita que
  se encuentra en Australia y Tasmania (Tremblay, 2006). En general las
  orquídeas epifitas tienen una historia de vida diferente a las
  terrestres, pero no se ha evaluado si estas diferencias son relevantes
  para la dinámica poblacional cuando uno considera el ambiente templado
  versus tropical.
\item
  La historia de vida de las especies terrestres en área templadas han
  sido bastante estudiando, pero en ambientes tropicales los estudios
  son escasos. Los que hemos encontrado solo son los de
  \emph{Oeceoclades maculata} (Riverón-Giró et al., 2019), \emph{Phaius
  australis} (Simmons, 2016) y \emph{Spathoglottis plicata} (Falcón,
  Ackerman \& Tremblay, 2017). Dado que el ambiente tropical es menos
  variable en temperatura, cual efecto si algunos tienen alhunos sobre
  la historia de vida de las orquídeas terrestres tropical comparado a
  las especies templadas.
\end{itemize}

\chapter{Preguntas evolutivas}\label{preguntas-evolutivas}

\begin{itemize}
\item
  ¿La historia de vida de los diferentes grupos de orquídea ha
  influenciado la diversidad taxonómica?
\item
  ¿Cuál es el impacto de la reproducción clonal sobre la dinámica
  poblacional y especialmente sobre los procesos evolutivos?
\item
  ¿La velocidad evolutiva es diferente entre una orquídea miniatura o de
  ramilla, dado que las miniatura tienen ciclos de vida más cortos, y
  por ende alcanzan una edad reproductiva en menor tiempo?
\end{itemize}

\chapter{Conclusión}\label{conclusiuxf3n}

Lo anterior son solamente una muestra de las preguntas que pueden ser
abordadas con los modelos de dinámica poblacional. Los autores esperamos
que este libro sea un aporte a la comprensión de la dinámica poblacional
de las orquídeas y que sirva como una guía para futuros estudios en este
fascinante grupo de plantas. La diversidad de preguntas que pueden ser
abordadas con los modelos de dinámica poblacional es amplia y variada, y
esperamos que este libro inspire a otros a explorar y responder estas
preguntas en el futuro.

\phantomsection\label{refs}
\begin{CSLReferences}{1}{0}
\bibitem[\citeproctext]{ref-ackerman1996seedling}
Ackerman JD, Sabat A, Zimmerman J. 1996. Seedling establishment in an
epiphytic orchid: An experimental study of seed limitation.
\emph{Oecologia} 106:192--198.

\bibitem[\citeproctext]{ref-ackerman2025persistent}
Ackerman JD, Tremblay R, Arias T, Zotz G, Sharma J, Salazar GA, Kaur J.
2025. Persistent habitat instability and patchiness, sexual attraction,
founder events, drift and selection: A recipe for rapid diversification
of orchids. \emph{Plants} 14:1193.

\bibitem[\citeproctext]{ref-batty2001constraints}
Batty A, Dixon K, Brundrett M, Sivasithamparam K. 2001. Constraints to
symbiotic germination of terrestrial orchid seed in a mediterranean
bushland. \emph{New {P}hytologist} 152:511--520.

\bibitem[\citeproctext]{ref-ehrlen2000dispersal}
Ehrlén J, Eriksson O. 2000. Dispersal limitation and patch occupancy in
forest herbs. \emph{Ecology} 81:1667--1674.

\bibitem[\citeproctext]{ref-falcon2017quantifying}
Falcón W, Ackerman JD, Tremblay R. 2017. Quantifying how acquired
interactions with native and invasive insects influence population
growth rates of a non-indigenous plant. \emph{Biological Invasions}
19:895--911. DOI:
\href{https://doi.org/10.1007/s10530-016-1318-8}{10.1007/s10530-016-1318-8}.

\bibitem[\citeproctext]{ref-garcia2010living}
Garcı́a MB, Goni D, Guzmán D. 2010. Living at the edge: Local versus
positional factors in the long-term population dynamics of an endangered
orchid. \emph{Conservation Biology} 24:1219--1229.

\bibitem[\citeproctext]{ref-light1998factors}
Light MH, MacConaill M. 1998. Factors affecting germinable seed yield in
\emph{{C}ypripedium calceolus var. {p}ubescens} ({W}illd.) {C}orrell and
\emph{{E}pipactis helleborine} ({L}.) {C}rantz ({O}rchidaceae).
\emph{Botanical Journal of the Linnean Society} 126:3--26.

\bibitem[\citeproctext]{ref-morris2020biotic}
Morris WF, Ehrlén J, Dahlgren JP, Loomis AK, Louthan AM. 2020. Biotic
and anthropogenic forces rival climatic/abiotic factors in determining
global plant population growth and fitness. \emph{Proceedings of the
National Academy of Sciences} 117:1107--1112.

\bibitem[\citeproctext]{ref-ramirez2021host}
Ramı́rez Martı́nez A. 2021. Host tree effect on demography and phenology
of epiphytic species. PhD thesis Thesis. Instituto Polit{é}cnico
Nacional. Centro Interdisciplinario de Investigaci{ó}n.

\bibitem[\citeproctext]{ref-riveron2019spatio}
Riverón-Giró FB, Raventós J, Damon A, Garcı́a-González A, Mújica E. 2019.
Spatio-temporal dynamics of the invasive orchid \emph{{O}eceoclades
{m}aculata} ({O}rchidaceae), in four different habitats in southeast
{C}hiapas, {M}exico. \emph{Biological Invasions} 21:1905--1919.

\bibitem[\citeproctext]{ref-simmons2016Viability}
Simmons C. 2016. Conservation genetics and population viability of the
endangered swamp orchid \emph{{P}haius {a}ustralis} ({O}rchidaceae) in a
changing world. PhD thesis Thesis. Faculty of Science, Health,
Education; Engineering, University of the Sunshine Coast, Queensland,
Australia: University of the Sunshine Coast.

\bibitem[\citeproctext]{ref-tremblay1992trends}
Tremblay Rl. 1992. Trends in the pollination ecology of the
{O}rchidaceae: Evolution and systematics. \emph{Canadian Journal of
Botany} 70:642--650.

\bibitem[\citeproctext]{ref-tremblay1994var}
Tremblay R. 1994. Frequency and consequences of multi-parental
pollinations in a population of \emph{{C}ypripedium {c}alceolus} var.
\text{pubescens} ({O}rchidaceae). \emph{Lindleyana} 9:161--167.

\bibitem[\citeproctext]{ref-tremblay2006effect}
Tremblay R. 2006. The effect of flower position on male and female
reproductive success in a deceptively pollinated tropical orchid.
\emph{Botanical Journal of the Linnean Society} 151:405--410.

\bibitem[\citeproctext]{ref-wilcock2002pollination}
Wilcock C, Neiland R. 2002. Pollination failure in plants: Why it
happens and when it matters. \emph{Trends in Plant Science} 7:270--277.

\end{CSLReferences}


\backmatter


\end{document}
